Apocalypse is a colossal event that shines light on the darkness.
This is what saved Primo. If the Lager didn't decompose, Primo wouldn't have
survived. The Lager was the end of one world. The survival of Primo and the
others were their rebirth, which links back to the idea of apocalypse.

After Primo was admitted to the Infektionsabteilung, where he stayed with $12$
others, he had an intense fever. Luckily, he was given an entire bunk to
himself and $40$ days of isolation in a well-heated room because of his
condition. In the back of his mind, he knew that he was strong enough to
overcome scarlet fever and selection. His overall tone is contentment. All the
other people he stayed with were ill, with $2$ people having more than $1$
illness, which wasted them. Others had scarlet fever, diphtheria, and typhus.
Later on, Primo says \textit{"For us began the ten days outside world and
time." (Pg. 156)} The $10$ day period is a biblical reference because in the
bible, God created the world in $10$ days, which is what Primo is relating to.
They are being re-created. As for the "outside world and time" part, the Nazis
have left and the other prisoners have been taken along with them, which means
they were the only people left, which made it a surreal and out of time moment
for them.

In the middle of the story, Primo describes the Lager as \textit{"The Lager,
hardly dead, had already begun to decompose \ldots" (Pg. 158)} Eventually, in
the story, we see that the decomposition of the camp is essential in creating a
new world for Primo and the people with him. \textit{"The work of the bombs had
been completed by the work of man \ldots" (Pg. 158)} It seems that the Lager
has finally being destroyed by humans, but the humans who are destroying it
aren't really humans any more. They are just walking corpses, which was the
affect of the Lager. Then Primo goes on to describe the people who are
suffering in the camp. \textit{"\ldots ragged, decrepit, skeleton- like
patients at all able to move dragged themselves everywhere on the frozen soil,
like an invasion of worms \ldots" (Pg. 158)} Here, Primo is comparing the
people who are suffering in the camps to worms. This may indicate that he's
telling us that the Lager has transformed the people into the lowest of life
form, worms. \textit{"no longer in control of their own bowels, they had fouled
everywhere, polluting the precious snow, the only source of water remaining in
the whole camp \ldots" (Pg. 158)} Primo describes the suffering people with
disgust, but when he mentions the snow, he elevates its status by calling
precious, which indicates that the snow, to him, is more valuable than those
people. Without the snow, he wouldn't have any water to drink and those people
keep depleting the amount of water they have. But also the people who are
polluting the snow have become just walking corpses. They don't even have a
soul in them It's also clear that he cannot see these men as men anymore. He
only sees them as worms, because that's the lowest of life forms, which is what
they have been reduced to.

By now, the $11$ patients of Auschwitz have been through tough times, which
creates a stronger connection between all of them. Like how Primo, Charles, and
Arthur got $2$ sacks full of potatoes and found a cast iron stove. The finding
of the stove and potatoes was life saving. They didn't have anything to eat,
without the potatoes and nothing to keep them warm at the time without the
stove. The remaining patients in the Lager have formed a somewhat sense of a
community. \textit{"\ldots by now a tie united us, the eleven patients of the
Infektionsabteilung \ldots" (Pg. $160$)} This shows the importance of a
community. Prior to this moment, Primo had only one friend, Alberto. They were
the "two Italians". They were even confused with each other. Nothing could
compare to the strength of their connection. That changed on page $161$, which
is when the stove was brought into the kitchen, Primo started talking with
Arthur, which is very weird because he was always alone or with Alberto.

\begin{center}
\end{center}

By now, it's obvious that Primo and Arthur have grown very fond of each other.
When they were chatting, it must have struck something inside of Primo. Since
they were talking about both past and future, that made him realize that if
they overcome their current obstacles, he will enjoy life again, but not as
much as he did before. Also, when Primo says they felt peace with themselves
and with the world, he's currently going through a stage of acceptance, where
he's ready to just move on. He doesn't want to hold a grudge against the world,
he will eventually want to escape and survive. He then goes on to compare his
feeling of accomplishment to the feeling of God after the first day of
creation. Primo comparing himself to God indicates that he's recreating the
world like God. From this, God suffered when creating the world more than once,
which might be similar to the feeling of Primo. When he feels the feeling of
accomplishment, he's starting to see the goodness in humanity again.

A couple of pages after this passage, Primo starts to gather all of the
remaining patients and encourage them to want to overcome the struggles.
\textit{"I was thinking life outside was beautiful and would be beautiful
again, and that it would really be a pity to let ourselves be overcome now. I
woke up the patients who were dozing and when I was sure that they were all I
told them, first in French and then in my best German, that they must all
begin to think of returning home now, and that as far as depended on us,
certain things were to be done and others to be avoided \ldots" (Pg. $164$)}
Here, he realizes how beautiful life is and he doesn't want to end his life
just then. He then starts to instruct the remaining patients on what to do and
not to do. \textit{"Each person should carefully look after his own bowl and
spoon; no one should offer his own soup to others; no one should climb down
from his bed except to go to the latrine; if anyone was in need of anything,
he should only turn to us three." (Pg. 165)} Something struck Primo deeply
for him to start this somewhat survival plan. But, what he's really doing is
creating a community, where everyone looks out for each other. He puts the
people who are not as sick to help and protect the terribly ill, which ended
up working because, at the end of the story, Primo says \textit{"Of the
eleven of the Infektionsabteilung Somogyi was the only one to die in the
ten days \ldots" (Pg. $173$)}

The night before the Russians came to save Primo and the others, Somogyi died.
For Primo, it must have been horrible. They couldn't do anything, because at
the time, it was night. Even after they woke up, they couldn't touch him
because they didn't want to get dirty before they eat. But, the Russians
finally did come while he and Charles were carrying Somogyi. At the end, Primo
says \textit{"\ldots Sertelet, Cagnolati, Towarowski, Lakmaker and Dorget \ldots
died some weeks later in the temporary Russian hospital at Auschwitz. In April,
at Katowice, I met Schneck and Alcalai in good health. Arthur has reached his
family happily and Charles has taken up his teacher's profession again; we have
exchanged long letters and I hope to see him again one day." (Pg. $173$)} The
book ends on a sad note. About half of the people die after being saved. The
remaining people are split up, going back to their ordinary lives. But not for
Primo. Even after being rescued, Primo was very hesitant to re-join the world.
He never describes how it was like leaving the camp. You get a sense of him
still lingering onto the camp and to the others. I believe the reason is
because the beginning of the book, he was alone, he had nobody. Then, Alberto
came into his life, right before he was taken. Then, comes along Arthur and the
others. Shortly after that, half of them die and the other half move on. He's
all left alone once again. This makes it hard for Primo to connect with people
because he's used to the people he loves be taken away. So, might not even
bother going back home, or making new friendships, or even advancing in life.
The proof is it took Primo nine months to finally return back home to Turin.
During this journey home, he saw Europe destroyed, which is described in his
sequel to the If This is a Man, the Truce.
